\solchapter{sr}

\sol{sr}{1} The trip takes $15\un{km}/(0.998c) = 50\un{\mu s}$. Without time dilation the survivors are $e^{-50/2.2}\approx10^{-10}$. With $\gamma = 15.8$ the lab-frame lifetime is $35\un{\mu s}$ and $e^{-50/35}\approx24\%$ survive. In the muon's frame, the atmosphere is contracted to $15/15.8 = 0.95\un{km}$ and rushes past in $3.2\un{\mu s}$ --- the same $1.4$ lifetimes, the same $24\%$.

\sol{sr}{2} $mc^2 = 10^{-3}\times9\times10^{16} = 9\times10^{13}\un{J}$, about $21$ kilotons of TNT. Fifteen kilotons, $6.3\times10^{13}\un{J}$, is about $0.7\un{g}$ of mass.

\sol{sr}{3} $u = (0.8+0.8)/(1+0.64)\,c = 0.976c$. Galileo would have said $1.6c$.

\sol{sr}{4} $v^2/2c^2 = (3.9\times10^3)^2/(1.8\times10^{17}) = 8.5\times10^{-11}$; times $86\,400\un{s}$, about $7\un{\mu s}$ per day. Light travels $2\un{km}$ in that time, so uncorrected position errors would grow by about $2\un{km}$ per day from this effect alone.

\sol{sr}{5} In the pion's rest frame the muon and neutrino have equal and opposite momenta. Energy conservation gives $E_\mu = (m_\pi^2+m_\mu^2)c^2/2m_\pi = (139.6^2+105.7^2)/279.2\un{MeV} = 109.8\un{MeV}$, so its kinetic energy is $4.1\un{MeV}$.

\sol{sr}{6} At threshold the four final particles are at rest in the centre-of-mass frame, so $s = (4m_pc^2)^2$. For a beam of energy $E$ on a proton at rest, $s = 2m_p^2c^4 + 2Em_pc^2$. Setting these equal, $E = 7m_pc^2$, and the kinetic energy is $6m_pc^2 = 5.6\un{GeV}$. Most of the energy is not available for making mass: momentum conservation forces the products to move forward, carrying kinetic energy.

\sol{sr}{7} $c^2t'^2-x'^2 = \gamma^2[(ct-\beta x)^2 - (x-\beta ct)^2] = \gamma^2(1-\beta^2)(c^2t^2-x^2) = c^2t^2-x^2$. Since $\gamma^2-(\gamma\beta)^2 = 1$, we can write $\gamma = \cosh\phi$ and $\gamma\beta = \sinh\phi$, with $\tanh\phi = \beta$. Two boosts multiply as hyperbolic rotations, and the hyperbolic addition formulas give $\phi = \phi_1+\phi_2$ --- which, written in terms of velocities, is exactly the velocity-addition formula.

\sol{sr}{8} Squaring $p_e' = p_\gamma + p_e - p_\gamma'$ and using $p_\gamma^2 = p_\gamma'^2 = 0$, $p_e^2 = p_e'^2 = m^2c^2$:
\[
  0 = p_\gamma\cdot p_e - p_\gamma\cdot p_\gamma' - p_e\cdot p_\gamma' = Em - \frac{EE'}{c^2}(1-\cos\theta) - E'm .
\]
Dividing by $mEE'$: $1/E' - 1/E = (1-\cos\theta)/mc^2$, which with $E = hc/\lambda$ is $\Delta\lambda = (h/m_ec)(1-\cos\theta)$.

\sol{sr}{9} The Earth twin waits $8/0.8 = 10$ years. The traveller's clock runs slow by $\gamma = 5/3$, recording $6$ years (in its own frame each leg is $4/\gamma = 2.4$ light-years long, crossed in $3$ years). The asymmetry: the Earth twin stays in one inertial frame, while the traveller changes frames at the turnaround. The turnaround is the event $(t,x) = (5\un{yr}, 4\un{ly})$. The outgoing frame's line of simultaneity through it, $t - 0.8x = $ const, meets the Earth's worldline at $t = 1.8$ years; the returning frame's line, $t + 0.8x = $ const, meets it at $t = 8.2$ years. So in the traveller's accounting, $6.4$ years of Earth time ``happen'' during the turnaround, and $1.8 + 6.4 + 1.8 = 10$. On the spacetime diagram, the Earth's straight worldline has the greatest proper time; the bent one has less.

\sol{sr}{10} ``Both doors shut at the same time'' is a statement about simultaneity, which is frame-dependent. In the pole's frame, the far door shuts and reopens \emph{before} the front of the pole reaches it, and the near door shuts only \emph{after} the back has passed it; at no moment is the pole inside with both doors shut. The two door events are spacelike separated, so their order depends on the frame, and both descriptions are correct. If the far door stays shut, the front of the pole stops; the back cannot learn of this faster than light, so it keeps coming, and the pole is crushed. There are no rigid bodies in relativity. Both frames agree on the wreckage.

\sol{sr}{11} Here $c^2/g = 0.97\un{ly}$ and $c/g = 0.97\un{yr}$. After $10$ years, $g\tau/c = 10.3$ and $x\approx0.97\cosh(10.3)\approx1.4\times10^4\un{ly}$. After $20$ years, $x\approx4\times10^8\un{ly}$ --- far beyond the Andromeda galaxy. To cover $13\,000\un{ly}$ to the midpoint: $\cosh(g\tau/c) = 1 + 13\,000/0.97$, so $g\tau/c = 10.2$ and $\tau\approx10$ years; with deceleration, about $20$ years of crew time (and $26\,000$ years on Earth). The catch is fuel: even a perfect photon rocket, with exhaust at $c$, needs a mass ratio $e^{\Delta\phi}$ for a total rapidity change $\Delta\phi$, here $e^{20}\approx10^9$. And at $\gamma\approx10^4$, interstellar hydrogen atoms hit the ship like cosmic rays.
